\documentclass[11pt,a4paper]{article}
\usepackage[T1]{fontenc}
\usepackage[utf8]{inputenc}
\usepackage{amsmath,amssymb,xcolor}
\usepackage[margin=25mm]{geometry}
\usepackage[colorlinks=true,urlcolor=blue]{hyperref}
\definecolor{nexorablue}{HTML}{2459B5}
\setlength{\parindent}{0pt}
\setlength{\parskip}{7pt}
\setlength{\emergencystretch}{2em}
\renewcommand{\familydefault}{\sfdefault}
\begin{document}

{\small\bfseries\color{nexorablue} NEXORA CPGE}\par
{\small MPSI / NEX-CPL-06 / Corrigé}\par
{\LARGE\bfseries Centre d’une composition de similitudes}\par
\vspace{0.25em}\hrule\vspace{0.6em}

\section*{Composée et point fixe}
On obtient $(T\circ S)(z)=2iz+3-i$. Le point fixe vérifie $(1-2i)z=3-i$, soit $z=1+i$ en utilisant le conjugué du dénominateur.

\section*{Interprétation}
En notant $c=1+i$, l’égalité $(T\circ S)(z)-c=2i(z-c)$ identifie une similitude de rapport $2$ et d’angle $\pi/2$ modulo $2\pi$.


\end{document}
